\documentclass[12pt]{exam}
\usepackage{amsmath,amsthm,amsbsy,amsfonts}
 
\boxedpoints
\pointsinmargin
\addpoints
 
\newcommand{\Z}{\mathbb Z}
\newcommand{\E}{\mathbb E}
\newcommand{\R}{\mathbb R}
\newcommand{\Q}{\mathbb Q}
 
\pagestyle{headandfoot}
%\extraheadheight{.4in}
%\extrafootheight{-1in}
%\extrawidth{.7in}
\firstpageheader{Calculus II}{\bf\Large Worksheet}{}
 
\begin{document}

\begin{questions}

\question If the amount of mayhem is constant in an episode of the Muppet Show, there is a simple
formula that gives the amount of mess created:
\[ \text{Mess = Mayhem} * \text{Time} \]
That is, the mess is the product of the mayhem and the length of time it lasts. Typically,
however, the mayhem varies over time. In a certain 30 minute episode the amount of
mayhem was given by the function $M(t) = 4 \ln(t + 1)$, where $0\leq t \leq 30$. Describe in
complete detail how integration can be used to find the total mess in the episode.
\vspace{.3in}

\question  It is a well-known fact that the anger generated by a fan
during a football game is caused by both the frustration they experience
and the length of time it occurs.  The anger produced by a 
 constant ammount of frustration is given by the formula
\[
\text{Anger = }\frac{1}{3} \left(\text {Frustration}\right)^2 * \text{Time}.
\]
In the first quarter of a game, the Frustration was given by
$\displaystyle F(t)= \sin\left(\frac{\pi}{30} t\right)  $, where $0\leq t\leq 15$.
\vspace{.2in}

\textbf{Describe in complete detail} how integration can be used to find
the total Anger in the first quarter.  Discuss the role played by Riemann 
sums, and set up an integral to find this Anger.
\vspace{.2in}

\question A conical tank  has a radius of 5 ft and a height of 13 ft.
Set up an integral to determine how much work is done in emptying it if
\begin{parts} 
\part It is filled with milk (which weighs 64.5 $lb/ft^3$).
\vspace{.3in}
\part It is filled to a depth of 9 ft and the milk must be pumped out
a 2 ft pipe at the top of the tank.
\vspace{.3in}
\end{parts}

\question  A dam has a triangular gate (triangle pointing upwards) with base 4 ft
and height 5 ft.
The top of the gate (the point of the triangle) 
is 3 feet below the water level.  Set up an integral
to find the force on
the gate. (Recall the density of water is 62.4 lb/ft$^3$.)
\vspace{.3in}

\question Set up integrals to find the mass, moment about the $x$-axis, and
moment about the $y$-axis of the lamina bounded by $y=\sqrt{x}$ and
$y=\displaystyle \frac{1}{3} x$.  
Give the formula to find the coordinates of the center of mass.
\vspace{.3in}

\question A spring has a natural length of 5 ft.  A force of
4 lb stretches this spring 1/4 ft.  
\begin{parts}
\part Find the work done in stretching this spring from its natural
length to 7 ft.
\vspace{.3in}
\part find the work done in stretching this spring from 8 ft to 9 ft.
\end{parts}
\vspace{.3in}

\question How much work is done in emptying a cylindrical tank of 
radius 5 ft and height 10 ft by pumping the the water through a 
pipe extending 4 ft above the tank.  (Hint: Water weighs 62.4 lb/$ft^3$.)
\vspace{.3in}

\question A rectangular gate built into a dam is 4 ft high and 8 ft wide.
The top of the gate is 5 ft below the surface of the water.  What is the 
force on the gate?
\vspace{.3in}

\question What would be  the force if the gate were the upper half of disc of
radius 4 (and the top of the gate is 5 ft below the surface of the water)?
\vspace{.3in}

\question Find the mass, moments, and center of mass of a the
triangular lamina bounded by the $x$- and $y$- axes and $y=4-2x$.
(As always, assume the density of this lamina to be $1$.)










\end{questions}
\end{document}